% Part: first-order-logic
% Chapter: axiomatic-deduction
% Section: deduction-theorem

% verification of properties of provability needed for maximally
% consistent sets in the completeness chapter.

\documentclass[../../../include/open-logic-section]{subfiles}

\begin{document}

\iftag{FOL}
      {\olfileid[zh]{fol}{axd}{ded}}
      {\olfileid[zh]{pl}{axd}{ded}}

\olsection{演绎定理}

正如我们所看到的，在公理系统中给出!!{derivation}颇为繁琐，而且!!{derivation}可能难以找到。与其真正写出长串的!!{formula}，通常更简单的做法是利用少数几个简单结果来论证这样的!!{derivation}存在。我们已经确立了三个这样的结果：
\olref[ptn]{prop:reflexivity}说，当我们知道 $!A \in
\Gamma$ 时，总可以断言 $\Gamma
\Proves !A$。\olref[ptn]{prop:monotonicity}说，如果 $\Gamma \Proves !A$，则亦有 $\Gamma \cup \{!B\} \Proves !A$。而
\olref[ptn]{prop:transitivity}蕴含：如果 $\Gamma \Proves !A$ 且
$!A \Proves !B$，则 $\Gamma \Proves !B$。这里还有另一个简单
结果，即分离规则的一个「元」版本：

\begin{prop}
  \ollabel{prop:mp} 如果 $\Gamma \Proves !A$ 且 $\Gamma \Proves !A \lif
  !B$，则 $\Gamma \Proves !B$。
\end{prop}

\begin{proof}
  我们有 $\{!A, !A \lif !B\} \Proves !B$：
  \begin{derivation}
    1. & $!A$ & 假设\\
    2. & $!A \lif !B$ & 假设\\
    3. & $!B$ & 1, 2, MP
  \end{derivation}
  由 \olref[ptn]{prop:transitivity}，$\Gamma \Proves !B$。
\end{proof}

在这一语境中我们将用到的最重要的结果是演绎定理：

\begin{thm}[演绎定理]
\ollabel{thm:deduction-thm} $\Gamma \cup \{!A\} \Proves !B$ 当且仅当 $\Gamma \Proves !A \lif !B$。
\end{thm}

\begin{proof}
「若」方向是直接的。如果 $\Gamma \Proves !A \lif !B$，则由
\olref[ptn]{prop:monotonicity} 亦有 $\Gamma \cup \{!A\} \Proves !A \lif !B$。再由
\olref[ptn]{prop:reflexivity}，$\Gamma \cup \{!A\} \Proves !A$。于是，由 \olref{prop:mp}，$\Gamma \cup
\{!A\} \Proves !B$。

对于「仅当」方向，我们以推演的长度进行归纳，其中!!{derivation}是 $!B$ 从 $\Gamma \cup \{!A\}$ 出发的推演。

对于归纳基始，我们证明该断言对所有长度为 $1$ 的!!{derivation}成立。$!B$ 从 $\Gamma \cup \{!A\}$ 出发的长度为 $1$ 的!!^a{derivation}恰由 $!B$ 本身构成；若它是正确的，则 $!B$ 要么 $\in \Gamma \cup \{!A\}$，要么是一条公理。如果 $!B \in \Gamma$ 或者是公理，则 $\Gamma \Proves !B$。我们亦有 $\Gamma
\Proves !B \lif (!A \lif !B)$，这来自 \olref[prp]{ax:lif1}，而
\olref{prop:mp} 给出 $\Gamma \Proves !A \lif !B$。如果 $!B \in \{ !A\}$，则 $\Gamma \Proves !A \lif !B$，因为最后的那个!!{sentence}~$!A
\lif !B$ 与 $!A \lif !A$ 相同，而我们已在
\olref[pro]{ex:identity} 中!!{derive}出了它。

对于归纳步骤，设 $!B$ 从
$\Gamma \cup \{!A\}$ 出发的!!a{derivation}以某一步骤$~!B$ 结尾，而该步骤由分离
规则证成。（如果它不是由分离规则证成的，则 $!B \in \Gamma$、$!B
\ident !A$，或 $!B$ 是一条公理，此时与归纳基始相同的推理适用。）那么在!!{derivation}中某两个较早的步骤是 $!C \lif !B$ 与 $!C$，对某个!!{formula}~$!C$
而言，即 $\Gamma \cup \{!A\} \Proves !C \lif !B$ 且 $\Gamma \cup \{!A\}
\Proves !C$，而相应的!!{derivation}更短，因此归纳
假设适用于它们。于是我们两者皆有：
\begin{align*}
  & \Gamma \Proves !A \lif (!C \lif !B); \\
  & \Gamma \Proves !A \lif !C.
\end{align*}
但亦由 \olref[prp]{ax:lif2}，有
\[
\Gamma \Proves (!A \lif (!C \lif !B)) \lif
((!A\lif !C)  \lif (!A \lif !B)),
\]
而两次应用 \olref{prop:mp} 即给出
$\Gamma \Proves !A \lif !B$，如所要求。
\end{proof}

注意 \olref[prp]{ax:lif1} 与 \olref[prp]{ax:lif2} 是如何被选定，
恰是为了使演绎定理成立。

下面是一些关于!!{derivability}的有用事实，我们把它们留作习题。

\begin{prop}
  \ollabel{prop:derivfacts}
\begin{enumerate}
\item $\Proves (!A \lif !B) \lif ((!B \lif !C)
  \lif (!A \lif !C)$； \ollabel{derivfacts:a}
\item 如果 $\Gamma \cup \{ \lnot !A\}
  \Proves \lnot !B$，则 $\Gamma \cup \{ !B\} \Proves
  !A$（换质位）； \ollabel{derivfacts:b}
\item  $\{ !A, \lnot!A\} \Proves
    !B$（爆炸原理）； \ollabel{derivfacts:c}
\item  $\{ \lnot\lnot!A\} \Proves
  !A$（双重否定消去）；\ollabel{derivfacts:d}
\item 如果 $\Gamma \Proves \lnot\lnot!A$，则 $\Gamma \Proves
  !A$；\ollabel{derivfacts:e}
\end{enumerate}
\end{prop}

\tagprob{FOL}
\begin{prob}
  证明 \olref[fol][axd][ded]{prop:derivfacts}
\end{prob}
\tagendprob

\tagprob{notFOL}
\begin{prob}
  证明 \olref[pl][axd][ded]{prop:derivfacts}
\end{prob}
\tagendprob

\end{document}
